%----------------------------------------------------------------
%### Loading Packages ###
%
%## Encoding, Typography & Language ##
\usepackage[utf8]{inputenc}
\usepackage[T1]{fontenc}
\usepackage{babel} 
\usepackage{lmodern,csquotes}
\PassOptionsToPackage{mathscr}{eucal}\usepackage{amsfonts,eucal,mathrsfs,textcomp} 
\usepackage{ulem}
\normalem
%\usepackage{newtxtext}

%## Colors, Boxes & Figures ##
\usepackage[usenames,x11names,table]{xcolor} %%colour names in 'xcolor-palettes.pdf' of the 'Templates_labOMC' repo. For printing, the [cmyk,dvipsnames] colour model is preferred.
\usepackage{fancybox} %%more framed boxes
\usepackage[breakable]{tcolorbox} %%coloured boxes with optional title
\usepackage{float} %%more control on any floating environments 
\usepackage{array,dcolumn,tabularx,multirow,longtable,booktabs} %%more control of 'tabular' and 'array' environments
\usepackage{graphicx,wrapfig,caption,sidecap} %% Compiling with 'pdfLaTeX' supports .pdf,.pdf_tex .png, .jpg and .eps(indirectly) files. %Use '\includeinkscape' for .pdf_tex files exported from Inkscape or convert directly to TiKZ code with the 'svg2tikz' extension.

%## Maths & Physics ##
\usepackage{amsmath,amsthm,amssymb}
%\usepackage{cancel,trfsigns} %%math simplifications & signs for transforms respectively
%\usepackage[all]{xy} %%nice math diagrams
\usepackage{siunitx} %scientific & complex numbers, measurement & uncertainty, see 'latex-SIunitx.pdf' in the 'Templates_labOMC' repo.
\usepackage{physics} %typesetting for mathematical physics, see 'latex-PhysicsPack.pdf' in the 'Templates_labOMC' repo.
\AtBeginDocument{\RenewCommandCopy\qty\SI} %to fix the command \qty used by both siunitx & physics packages, favoring the former.
\usepackage{pgfplots} %graphs of symbolic functions and equations
%\usepackage{tikz-network} %%crystal lattice & complex networks drawing
%\usepackage{tikz-optics} %%optical drawing
%\usepackage{pstricks,pst-optexp,pst-optic} %%neat PSTricks optical drawing
%\usepackage{tikzorbital,chemfig,chemmacros,mhchem} %%chemical drawings/diagrams and formulae & safety
%\usepackage{pstricks,pst-labo} %%laboratory chemical glassware
%\usepackage{pstricks,dsptricks,pst-sigsys} %%signal processing plots/drawings
%\usepackage[RPvoltages,americanvoltages,americancurrents]{circuitikz} %%electrical circuit drawing with Rising Potential voltages
%\usetikzlibrary{babel} %%needed for electrical circuit drawing when running babel-french
%\usepackage{yquant} %%quantum circuit drawing

%## Miscellaneous Nice Tools ##
\usepackage{ccicons} %%Creative Commons icons, reusable content at <https://search.creativecommons.org/>
%\usepackage[draft]{showkeys} %%uncomment to display the label of each '\ref' key in the document output
%\setlength{\marginparwidth}{2cm} %%uncomment to force a temporary resolution of 'todonotes' margin warning
\usepackage{lipsum,todonotes} %%'\lipsum[_a-b_][_c-d_]' for dummy text paragraph #_a-b_, line #_c-d_ ; standard margin notes with '\marginpar{}' & colored ones with '\todo{}'
%\usepackage{media9,animate,tikz} %%tools for multimedia and general drawing
%\usepackage{pythontex} %%runs Python code in the LaTeX source files and typeset the output

%## Bibliography ##
%\usepackage[backend=biber,sorting=none]{biblatex} %%see 'cheatsheet-BibLaTeX.tex' in the 'Templates_labOMC' repo.
%\addbibresource{} %%import your .bib file in BibLaTeX format from Zotero

%## Hyperlinks ##
\usepackage{hyperref,xurl} %%'\href{_URL_}{_hypertext_}' to link the hypertext & '\url{_URL_}' to link the displayed web address, each internal '~\ref{_label_}' is also hyperlinked ; xurl for more line breaks.
\usepackage{doi} %%automatically resolves object identifiers such as IDs from doi, arXiv, OCLC, etc., use '\doi{}' if the former bugs.

%## Layout ##
%%Watch for conflicts with the overall document class.
\usepackage{geometry} %%page layout, see configuration below
\usepackage{fancyhdr} %%headers & footers, see configuration below
\usepackage{setspace} %%interline spacing tools ; unit relative to current font: em ~width of an 'M' (uppercase), typically equal to the font size in pt
\usepackage{enumitem} %%tools to change list format
%\usepackage{tasks,epigraph,boxedminipage,multicol,lscape,rotating,makeidx,showidx,glossaries}

%## Subdividing a document ##
\usepackage{import} %%\import{_path_}{_filename_} includes the desired file.
%\usepackage{subfiles} %also allows standalone compilation.
%\newcommand{\onlyinsubfile}[1]{#1} %%for content compiled only in the subfile.
%\graphicspath{{folder/}{../folder/}{folder/}{../folder/}} %The relative path to images, starting from each compilation file call, must be listed here otherwise the \includegraphics command will not work. The ../ stipulate one level up in the file tree.
%\usepackage[active,tightpage]{preview} %output cropped just around the content. For the whole document, option active->inactive.
%\setlength\PreviewBorder{20pt}

%---------------------------------------------------------
%## Global settings for packages ##
\sisetup{separate-uncertainty=true} %%for SIunitx package
\hypersetup{   %%for hyperref package
    breaklinks=true,
    linktocpage=true,
    colorlinks=true,
    urlcolor=blue,
    linkcolor=blue,
    citecolor=blue,
    filecolor=blue}
\pgfplotsset{compat=1.18,width=5cm} %%for pgfplots package, backwards compatibility with 'compat'
\usetikzlibrary{arrows} %%for PGF/TikZ

%---------------------------------------------------------
%### Defining command shortcuts & macros ###
%
\newcommand{\be}{\begin{equation}} %%for numbered equations, use instead \[...\] to display without numbering or $...$ to include directly inline.
\newcommand{\ee}{\end{equation}}
\newcommand{\bea}{\begin{align}} %%number each equation and align them on &= the equal sign and more generally on the ampersand & character. 
\newcommand{\eea}{\end{align}}
\newenvironment{eqsplit}{\equation\aligned}{\endaligned\endequation}
\newcommand{\bes}{\begin{eqsplit}} %%defined the same way as the align environment, but considered as one equation on multiple lines with one number overall.
\newcommand{\ees}{\end{eqsplit}}
\newcommand{\beg}{\begin{gather}} %%gathering equations one below another without aligning.
\newcommand{\eeg}{\end{gather}}
\DeclareMathOperator{\TLs}{\mathscr{L}} %%Laplace transform symbol \TLs
\newcommand{\TLd}[1]{\TLs\!\left[#1\right]} %%Laplace transform operator \TLd{}
\DeclareMathOperator{\TLsi}{\mathscr{L}^{-1}} %%inverse Laplace transform symbol \TLsi
\newcommand{\TLi}[1]{\TLsi\!\left[#1\right]} %%inverse Laplace transform operator \TLi{}
\DeclareMathOperator{\TFsi}{\mathscr{F}^{-1}} %%inverse Fourier transform \TFsi
\newcommand{\TFi}[1]{\TFsi\!\left[#1\right]} %%inverse Fourier transform operator \TFi{}
\DeclareMathOperator{\TFs}{\mathscr{F}} %%Fourier transform \TFs
\newcommand{\TFd}[1]{\TFs\!\left[#1\right]} %%Fourier transform operator \TFd{}
\renewcommand*{\Re}{\mathop{}\!\mathfrak{Re}} %%Real part \Re redefined with 2 characters fraktur typesetting
\renewcommand*{\Im}{\mathop{}\!\mathfrak{Im}} %%Imaginary part \Im redefined with 2 characters fraktur typesetting